\section{Experiments}

\subsection{Numerical Aperture estimation}
Numerical aperture of an optical fibre can be calculated via the basic trigonometry if one can measure the output beam diameter $D$ at a given length $L$ from the fibre's aperture.
The ratio of the beam radius to the mentioned length is a tangent of the acceptance angle.
Based on this, we may transftom the~\autoref{eqn:NA} into:
\begin{equation}
    NA = \sin\arctan{\frac{D}{2L}} = \frac{D}{2\sqrt{L^2 + \frac{D^2}{4}}}
\end{equation}

For the NA estimation a He-Ne laser operating at \SI{632.8}{nm} coupled into five different fibres was used.
The setup diagram is presented in \autoref{fig:setup_NA}.
\begin{figure}
    \centering
    \includegraphics[width=0.7\textwidth]{setup_NA.pdf}
    \caption{Experimental setup for the NA measurements.
      He-Ne laser (1) generates the light beam at \SI{632.8}{nm}, which is focused by a lens (2), so it can be propagated through the fibre (3) and finally be projected onto the screen (4)}
    \label{fig:setup_NA}
\end{figure}
The end of the fibre aiming the lens was mounted on a table with the ability to move in the XYZ axes, so the fibre could be aligned with the laser beam.
The screen was positioned at distances $L_n$ from the fibre aperture, so we could take an averaged value\footnote{Of the components, not the results!}.
The beam diameter was judged by eye\footnote{To obtain a precise measurement, diameter of the beam should be determined for the $I=\frac{I_{max}}{e^2}$.}.

% The commented content comes from the remote classes in 2022.
% Consider useful information and delete the unnecessary content.
% The fibres used in the experiments were:
% \begin{enumerate}
%     \item High-power applications optical fibre with $\phi_{core}=\SI{105}{\micro\metre}$.
%     \item Standard communication multi-mode optical fibre with $\phi_{core}=\SI{62.5}{\micro\metre}$.
%     \item Single-mode telecommunication fibre for $\lambda=\SI{1.5}{\micro\metre}$ with $\phi_{core}=\SI{10}{\micro\metre}$.
%     \item Some mysterious fibre.
%     \item Single-mode fibre for propagation of the pumping laser radiation at $\lambda=\SI{980}{\nano\metre}$.
% \end{enumerate}
% \begin{table}[]
%     \centering
%     \begin{tabular}{c||c|c|c||c}
%         Fibre no. & $\phi_{core} [\SI{}{\micro\metre}]$ & $L [\SI{}{\centi\metre}]$ & $D [\SI{}{\milli\metre}]$ & NA\\
%         \hline\hline
%         \multirow{4}{*}{1} & \multirow{4}{*}{105} & 10 & 19 & \multirow{4}{*}{0.084}\\
%         & & 15 & 24 \\
%         & & 20 & 33 \\
%         & & 25 & 42 \\
%         \hline
%         \multirow{2}{*}{2} & \multirow{2}{*}{62.5} & 4 & 2 & \multirow{2}{*}{0.025}\\
%         & & 8 & 4 \\
%         \hline
%         \multirow{2}{*}{3} & \multirow{2}{*}{10} & 5 & 5 & \multirow{2}{*}{0.040}\\
%         & & 10 & 7 \\
%         \hline
%         4 & -- & -- & -- & -- \\
%         \hline
%         \multirow{2}{*}{5} & \multirow{2}{*}{--} & 10 & 16 & \multirow{2}{*}{0.075}\\
%         & & 15 & 24 \\
%         & & 20 & 31 \\
%         & & 25 & 35 \\
%     \end{tabular}
%     \caption{$\phi_{core}$ -- fibre core diameter; $L$ -- distance from the fibre aperture to the screen; $D$ -- projected beam diameter}
%     \label{tab:NA}
% \end{table}

\begin{table}[]
    \centering
    \begin{tabular}{c||c|c||c}
        Fibre no. & $L [\SI{}{\centi\metre}]$ & $D [\SI{}{\milli\metre}]$ & NA\\
        \hline\hline
        \multirow{4}{*}{1} & 10 & 19 & \multirow{4}{*}{0.084}\\
        & 15 & 24 \\
        & 20 & 33 \\
        & 25 & 42 \\
        \hline
        \multirow{2}{*}{2} & 4 & 2 & \multirow{2}{*}{0.025}\\
        & 8 & 4 \\
        \hline
        \multirow{2}{*}{3} & 5 & 5 & \multirow{2}{*}{0.040}\\
        & 10 & 7 \\
        \hline
        4 & -- & -- & -- \\
        \hline
        \multirow{4}{*}{5} & 10 & 16 & \multirow{4}{*}{0.075}\\
        & 15 & 24 \\
        & 20 & 31 \\
        & 25 & 35 \\
    \end{tabular}
    \caption{$L$ -- distance from the fibre aperture to the screen; $D$ -- projected beam diameter}
    \label{tab:NA}
\end{table}

% Each fibre projected somehow a unique image: for (1) the projected pattern resembled a multi-mode pattern (of definitely high order); the "modes" in (2) were a little bigger than those in (1); for (3) we could observe just couple of modes; for (4) we could observe four distinguishable modes; with the last (5) fibre only one distinguishable mode was clearly visible.

The results of the measurements together with the NA estimation are shown in Table~\ref{tab:NA}.
% All but (3) oscillate around $NA=0.1$, and tweaking with the $D$ in the range of 30\% resulted in minor final value changes.

% All of these values were compared with the fibres offered by ThorLabs:

% (1) may be compared to the FG105LVA~\cite{FG105LVA}, a high-power multi-mode fiber with $\phi_{core}=\SI{105}{\micro\metre}$ and $NA=0.1\pm 0.015$.

% (2) may compared only with the GIF625~\cite{GIF625}, but its $NA=0.275$.

% (3) may be compared to the SM600~\cite{SM600}, a single-mode fiber with $NA=\numrange{0.10}{0.14}$.

% (4) was some mysterious fibre able to guide 4 distinguishable modes. As mentioned in Section~\ref{normalized_frequency}, $V$ required for such an action is $\frac{V^2}{2}>4$ for step-index and $\frac{V^2}{4}>4$ for graded-index fibre. The profile of this fibre seemed to be sharp (distinguishable high-intensity centre and profile's "tail"), so I would assume graded-index fibre, unfortunately ThorLabs does not provide information about the refractive indices of the core and cladding, therefore precise comparison is not possible.

% (5) for this fibre the normalized frequency had to be around 2.405, as sometimes second mode was visible, and assuming average value of square root of Eq.~\ref{eqn:v}, we can estimate that its $\phi_{core}$:
% \begin{equation*}
%     2.405 \approx \frac{\pi \phi_{core }}{\SI{632.8}{\nano\metre}}0.1\rightarrow\phi_{core}\approx\SI{4.8}{\micro\metre}.
% \end{equation*}

\subsection{Losses introduced by the fibre and its connectors}

For the loss' measurements we used the standard telecommunication single-mode fibres SMF28 terminated with the FC/APC connectors.
Each of the fibres' ends were inspected via the microscope for any dirtiness, and cleaned via rubbing its tip with a sheet of paper if necessary. 
The results are presented in~\autoref{tab:losses}.
\begin{table}
    \centering
    \begin{tabular}{c|c|c}
        Number of patchcords & Power [$\qty{}{\dBm}$] & Loss [$\qty{}{\dB}$]\\
        \hline\hline
        1 & 6.73 & 0.00 \\
        2 & 6.54 & 0.19 \\
        3 & 6.28 & 0.45 \\
        4 & 6.24 & 0.49 \\
        5 & 5.87 & 0.86
    \end{tabular}
    \caption{Losses measurements}
    \label{tab:losses}
\end{table}

The interesting measurement is change from 3 to 4 patchords with only \qty{0.04}{\dB} loss.
Such good value is probably the result of improving other connection during displacing the patcords.
Tightening the FC/APC connectors might decreate the signal if two angled faces of the optical fibre meet and push againts each other, bending them and thus displacing the cores.

The last measurement was a 24 km fibre spool.
Connected in one direction it had 1.43 dBm.
In the opposite direction it had 1.55 dBm, and after fiddling with the connectors it reached 1.87 dBm.
Altough one can elaborate on in-fibre phenomena, it was most probably, as before, the result of introducing core position changes due to poor connections.
If we would like to perform any on the above measurements correct, we would have to reconnect the fibres multiple times.

% Why tightening the FC/APC connector may decrease the signal?
% Because two angled faces of the optical fibre may meet and push againts each other, bending them and displacing the cores.

% Szpula - 24 km

% 1.43 dBm
% w drugą stronę 1.55 dBm, jak sie pokreci to 1.87 dBm

% \section{Conclusions}
% -- almost single mode for our red right (for sure has the smallest core diameter). Signle-mode for 980nm

% if a shorter wavelength propagate in the fiber it may be a multi mode

% Different diameter of the core -- different amount of the guided mode
% \footnote{Note, that the communication OFs are suited to use with the infrared light.}
